% Sumber: docs/EXPERIMENT.md (EXP-005/006/007), docs/EXP-007_ARCFACE_EVALUATION.md,
% docs/RESEARCH_FINDINGS.md, docs/R9_FINAL_STATISTICAL_VALIDATION.md, docs/RESEARCH_FRAMING.md,
% docs/proposal/validation/PROPOSAL_GUIDELINE_AUDIT.md, RESEARCH_DECISIONS.md (Keputusan 1, 2, 4).
\chapter{Pendahuluan}
\label{chap:introduction}

\section{Latar Belakang}
\label{sec:background}

Model pengenalan ekspresi wajah (\textit{facial emotion recognition}, FER) biasanya dibangun dan diuji hanya pada segelintir dataset acuan (FER2013, AffectNet, RAF-DB) yang masing-masing dikumpulkan dengan caranya sendiri: kamera tertentu, populasi crowd-sourcing tertentu, instruksi pelabelan tertentu \citep{li2020deep}. Apakah model yang dilatih dengan data semacam itu akan berperilaku sama ketika dihadapkan pada populasi yang tidak ada dalam dataset tersebut adalah pertanyaan empiris, bukan sesuatu yang bisa dijawab hanya dari angka akurasi di atas kertas. \citet{li2020deeper} menunjukkan hal ini secara langsung untuk dataset FER, dan \citet{buolamwini2018gender} menunjukkan, untuk tugas yang berkerabat yaitu verifikasi wajah, bahwa akurasi rata-rata bisa menyembunyikan perbedaan besar antar-subkelompok pada sistem komersial.

Tesis ini menguji pertanyaan tersebut pada satu model FER yang sudah jadi dan sudah dipakai di dunia nyata, yaitu HSEmotion \citep{savchenko2022hsemotion}, sebuah klasifier ringan terlatih yang didistribusikan untuk keperluan produksi, diterapkan pada sampel ekspresi wajah yang diambil dari satu film dokumenter tentang masyarakat Papua. Sampelnya kecil (227 potongan wajah), berasal dari satu sumber saja, dan tidak dikumpulkan khusus untuk tujuan ini; sampel ini tidak mendukung klaim apa pun tentang wajah masyarakat Papua secara umum, dan tesis ini tidak membuat klaim semacam itu. Yang bisa didukung oleh data ini adalah studi kasus yang terbatas tapi terukur dengan baik: mengukur di mana kesalahan model ini terkonsentrasi pada sampel ini, apakah sebuah langkah praproses gambar mengubah hal itu, dan apakah pendekatan ekstraksi fitur yang berbeda, yaitu klasifier Regresi Logistik di atas \textit{embedding} wajah ArcFace yang dibekukan \citep{deng2019arcface}, menghasilkan pola kesalahan yang terukur berbeda pada data yang sama.

\textit{Embedding} di sini maksudnya adalah representasi numerik suatu wajah: sebuah vektor angka yang merangkum ciri wajah tersebut, bukan gambar itu sendiri. "Dibekukan" (\textit{frozen}) berarti angka-angka dalam model ArcFace tidak diubah sama sekali; model tersebut hanya dipakai apa adanya untuk menghasilkan vektor tadi, lalu vektor itu yang dipakai sebagai masukan bagi klasifier baru yang justru dilatih dari nol menggunakan label tesis ini sendiri. Perbedaan antara "memakai model yang sudah ada" dan "melatih model baru" ini dijelaskan lebih rinci pada Bab~\ref{chap:methodology}, karena perbedaan ini menentukan apa yang sebenarnya dibandingkan dalam tesis ini.

Proyek di balik tesis ini juga menghasilkan satu jalur kerja lain yang terpisah: uji coba anotasi titik-fitur wajah (\textit{facial landmark}, yaitu titik-titik acuan seperti sudut mata, ujung hidung, garis bibir) pada kumpulan potongan wajah baru dari film sumber yang sama (disebut ``Papuan V2''). Jalur ini didokumentasikan dalam tesis ini sebatas tahap yang benar-benar sudah dicapai, yaitu usulan titik-fitur otomatis yang menunggu pemeriksaan manusia, dan diperlakukan di seluruh tesis ini sebagai pengembangan infrastruktur pendukung, bukan sebagai bukti yang menjawab pertanyaan FER di atas. Penulis telah menegaskan pembagian ini secara langsung (\texttt{RESEARCH\_DECISIONS.md}, Keputusan 1 dan 2): klasifikasi emosi adalah subjek utama tesis ini; uji coba titik-fitur adalah jalur paralel yang masih berjalan dan mungkin berguna untuk kerja lanjutan, bukan sekarang.

\section{Rumusan Masalah}
\label{sec:problem}

Catatan internal proyek ini awalnya membingkai pekerjaan ini sebagai ``FER untuk wajah berkulit gelap menggunakan ArcFace.'' Tinjauan internal menandai pembingkaian itu sebagai keliru secara metodologis: kalimat itu sudah menetapkan teknik (ArcFace) dan deskriptor populasi (``berkulit gelap'') sebelum ada bukti apa pun yang menunjukkan bahwa keduanya memang sudut pandang yang tepat. Perhitungan warna kulit yang sebenarnya dipakai dalam proyek ini adalah sebuah kelompok (\textit{bucket}) kecerahan piksel pada ruang warna LAB, dihitung langsung dari piksel potongan wajah, bukan ukuran pigmentasi kulit, etnis, atau kategori demografis yang tervalidasi. Sebaran terukur pada sampel ini sendiri (kira-kira 59\% Gelap, \textit{Dark}, berbanding 41\% Gelap-Sedang, \textit{Medium-Dark}, pada skala 227 potongan) bukanlah populasi seragam seperti yang diasumsikan oleh pembingkaian awal tadi.

Yang benar-benar ditunjukkan oleh bukti justru lebih sempit dan lebih konkret. Keluaran HSEmotion (dipakai apa adanya, tanpa dilatih ulang) pada sampel ini mencapai akurasi keseluruhan 22,47\%, jauh di bawah yang biasanya dilaporkan untuk model ini pada dataset latihnya sendiri, dengan satu pola kesalahan yang spesifik dan berulang: delapan belas dari tujuh puluh enam potongan berlabel-asli Netral diprediksi sebagai Marah (23,7\%). Pola ini muncul identik baik dihitung pada sampel penuh 227 potongan maupun pada subset 135 potongan yang dipakai belakangan untuk perbandingan klasifier. Sampel yang sama juga menunjukkan selisih akurasi yang cukup besar antara dua kelompok kecerahan yang dikarakterisasi: 44,0\% untuk potongan Gelap berbanding 14,3\% untuk potongan Gelap-Sedang. Apakah selisih ini berasal dari kelompok kecerahan itu sendiri, dari faktor tersembunyi dalam cara rekaman sumber diambil dan diterangi, atau dari komposisi ekspresi yang kebetulan berbeda pada tiap kelompok, bukan sesuatu yang bisa dipastikan dari satu sampel observasional sebesar ini, dan tesis ini tidak mencoba memastikannya. Selisih ini diperlakukan sebagai asosiasi terukur yang layak dilaporkan secara presisi, bukan sebagai bukti mekanisme sebab-akibat.

\section{Tujuan Penelitian}
\label{sec:objectives}

Tiga tujuan berikut mengikuti langsung dari tiga eksperimen yang sudah selesai, dan masing-masing berpadanan dengan pertanyaan penelitian pada \S\ref{sec:rq}:

\begin{enumerate}
  \item Mengkarakterisasi perilaku klasifikasi HSEmotion pada sampel 227 potongan, yaitu akurasi keseluruhan, struktur kesalahan klasifikasi, dan selisih akurasi antar-kelompok kecerahan, sebagai garis dasar (\textit{baseline}) untuk dua tujuan berikutnya.
  \item Menentukan apakah konfigurasi CLAHE adaptif yang sudah diimplementasikan dan diuji dalam proyek ini mengubah akurasi atau pola kesalahan garis dasar tersebut, dan melaporkan hasilnya apa adanya: satu kali uji coba eksploratif, bukan ablasi yang tervalidasi secara statistik.
  \item Menentukan apakah klasifier Regresi Logistik yang dilatih di atas \textit{embedding} ArcFace menghasilkan akurasi atau sebaran kesalahan yang terukur berbeda dari HSEmotion pada subset evaluasi bersama berjumlah 135 potongan, dilaporkan sebagai analisis retrospektif atas bukti yang sudah ada sebelum pertanyaan penelitian ini dirumuskan secara formal, bukan sebagai uji konfirmatif yang dirancang di muka (\S\ref{sec:rq}).
\end{enumerate}

Satu jalur kerja keempat yang sifatnya eksplisit sekunder, yaitu mendokumentasikan uji coba anotasi titik-fitur wajah Papuan V2 sejauh tahap yang sudah dicapai (pembuatan usulan otomatis, menunggu pemeriksaan manusia), berjalan paralel dan dilaporkan pada Bab~\ref{chap:methodology} \S\ref{sec:methodology-landmarks}. Jalur ini bukan tujuan dari penelitian FER yang dilaporkan pada Bab~\ref{chap:results} sampai \ref{chap:conclusion}.

\section{Ruang Lingkup dan Batasan Penelitian}
\label{sec:scope}

Semua yang diukur dalam tesis ini berasal dari satu film dokumenter berdurasi 96 menit. Sampel 227 potongan dan subset 135 potongan yang diturunkan darinya (\S\ref{sec:dataset-overview}) sama-sama berasal dari satu sumber ini, di bawah satu kondisi pencahayaan dan kamera, dan dilabel oleh satu orang anotator. Tidak satu pun dari tiga eksperimen memakai partisi uji terpisah dalam arti konvensional: EXP-005/006 melaporkan metrik langsung atas seluruh sampel 227 potongan, dan klasifier pada EXP-007 divalidasi dengan validasi silang berkelompok (\textit{grouped cross-validation}), bukan pembagian latih/uji yang tetap, karena tidak ada kolom identitas subjek yang bisa dipakai untuk membagi data secara aman (Bab~\ref{chap:methodology} \S\ref{sec:leakage}).

Batasan-batasan berikut berlaku untuk seluruh tesis ini dan dinyatakan sekali di sini, tidak diulang kata demi kata di setiap bab: kebenaran lapangan (\textit{ground truth}) dibuat oleh satu anotator saja, tanpa pendapat kedua yang tercatat; asal-usul rekaman sumber serta dasar persetujuan penggunaan wajah orang-orang di dalamnya tidak terdokumentasi dan berada di luar kemampuan tesis ini untuk diselesaikan; kelas Marah dan Jijik terlalu jarang munculnya (masing-masing satu dan dua kejadian dari 227) sehingga metrik per-kelas untuk keduanya tidak bermakna, dan inilah sebabnya setiap perbandingan formal pada Bab~\ref{chap:results} memakai reduksi ke lima kelas; dan subset 135 potongan adalah hasil penyaringan ketat dari sampel 227 potongan, bukan populasi kedua yang dikumpulkan secara independen, sehingga hasil yang berulang pada keduanya adalah pengulangan dalam satu sampel, bukan replikasi lintas sampel.

\section{Pertanyaan Penelitian}
\label{sec:rq}

\textbf{Struktur ini adalah arah yang diusulkan penulis untuk tesis ini, disetujui secara internal per revisi ini; belum ditinjau atau disetujui oleh dosen pembimbing (Pak Fadhil).} Struktur ini menggantikan rumusan satu kalimat sebelumnya yang berisiko menyiratkan bahwa ketiga eksperimen berbagi satu populasi evaluasi yang sama; struktur di bawah ini menjaga agar populasi 227 potongan dan 135 potongan tetap dipisahkan secara eksplisit (lihat \texttt{RESEARCH\_DECISIONS.md}, Keputusan 4, untuk catatan lengkap revisi ini).

\textbf{Pertanyaan Penelitian Utama.} Bagaimana kinerja model pengenalan ekspresi wajah terlatih yang bersifat umum, dilihat dari kesalahan klasifikasi dan perbedaan antar-subkelompok kecerahan kulit (sebagai proksi), pada sampel gambar ekspresi wajah dari satu film sumber, dan bagaimana konfigurasi praproses kontras serta klasifier alternatif yang diuji memengaruhi hasil yang diamati?

\begin{description}
  \item[RQ1: Karakterisasi garis dasar.] Berapa akurasi keseluruhan, pola kesalahan klasifikasi yang dominan, dan selisih akurasi antar-subkelompok kecerahan kulit (proksi) dari HSEmotion pada sampel 227 gambar?
  \item[RQ2: Praproses kontras.] Bagaimana konfigurasi CLAHE adaptif yang diuji memengaruhi akurasi klasifikasi dan pola kesalahan HSEmotion pada sampel 227 gambar yang sama?
  \item[RQ3: Klasifier alternatif.] Pada subset evaluasi lima-kelas berjumlah 135 gambar, bagaimana akurasi dan pola kesalahan Regresi Logistik berbasis \textit{embedding} ArcFace dibandingkan dengan HSEmotion?
\end{description}

Ada satu catatan integritas penelitian yang perlu ditulis secara eksplisit, bukan dibiarkan tersirat: perbandingan yang mendasari RQ3 (EXP-007) sudah dijalankan sebelum struktur pertanyaan penelitian ini ditulis dalam bentuk apa pun. Bab~\ref{chap:results} memperlakukan RQ3 sesuai dengan kenyataan ini: sebagai interpretasi atas bukti yang sudah ada, di mana pertanyaannya diajukan secara retrospektif terhadap hasil yang sudah ada, bukan sebagai uji yang dirancang di muka berdasarkan hipotesis untuk menjawabnya. Perbedaan ini memengaruhi seberapa kuat hasil tersebut bisa dipercaya: perbandingan yang didaftarkan di muka (\textit{prospectively registered}) dan analisis retrospektif atas data yang sama tidak punya bobot bukti yang sama, dan tesis ini tidak mengklaim yang pertama untuk RQ3.

\section{Kontribusi Penelitian}
\label{sec:contributions}

Kontribusi dari pekerjaan yang sudah selesai adalah gambaran yang terukur dan presisi tentang bagaimana satu model FER yang sudah dipakai dan satu alternatif berbasis representasi wajah berperilaku pada satu sampel yang berada di luar distribusi data latihnya, bukan klaim umum tentang ketahanan FER, keadilan (\textit{fairness}), atau wajah masyarakat Papua secara umum. Secara konkret: karakterisasi yang dapat direproduksi atas pola kesalahan dominan HSEmotion dan selisih akurasi antar-kelompok kecerahan kulit pada 227 potongan (RQ1); hasil satu kali uji coba yang dilaporkan dengan hati-hati bahwa konfigurasi CLAHE yang diuji tidak memperbaiki, bahkan sedikit memperburuk, akurasi HSEmotion pada sampel ini (RQ2); dan perbandingan teruji secara statistik namun tidak signifikan antara HSEmotion dan alternatif berbasis ArcFace pada 135 potongan, dilaporkan dengan catatan retrospektif di atas dan tanpa menafsirkan hasil tidak signifikan tersebut sebagai bukti bahwa kedua model setara (RQ3).

Uji coba titik-fitur wajah Papuan V2 adalah kontribusi terpisah yang masih berjalan, bagian dari upaya pengembangan dataset yang lebih luas dalam proyek ini. Sebagaimana didokumentasikan pada Bab~\ref{chap:methodology} \S\ref{sec:methodology-landmarks}, uji coba ini saat ini baru terdiri dari usulan titik-fitur yang dibuat otomatis dan menunggu pemeriksaan manusia; tidak ada klaim bahwa uji coba ini sudah mencapai tahap yang memengaruhi, atau bisa ditunjukkan memengaruhi, kinerja klasifikasi emosi.

\section{Sistematika Penulisan}
\label{sec:organization}

Bab~\ref{chap:literature-review} membahas literatur yang diperlukan untuk menempatkan HSEmotion, ArcFace, CLAHE, dan evaluasi bertingkat-subkelompok dalam konteks tiga pertanyaan penelitian tesis ini. Bab~\ref{chap:methodology} menjelaskan \textit{pipeline} (rangkaian langkah pemrosesan data), dua populasi dataset dan bagaimana satu diturunkan dari yang lain, serta alur kerja anotasi dan uji coba titik-fitur, dengan detail implementasi yang cukup agar hasil dapat direproduksi dari artefak yang ada. Bab~\ref{chap:results} menjawab RQ1, RQ2, dan RQ3 secara berurutan, masing-masing terhadap populasinya sendiri dan masing-masing dengan penjelasan eksplisit tentang apa yang didukung dan apa yang tidak didukung oleh hasil tersebut. Bab~\ref{chap:conclusion} ditutup dengan apa yang sudah dan belum bisa dijawab dari pertanyaan penelitian, serta apa yang masih terbuka untuk ditinjau dosen pembimbing dan untuk kerja lanjutan.
